\chapter{Grafy Reeba i funkcje Morse'a na powierzchniach}
\label{ch:tg-reeb}

\section{Od poziomic funkcji do grafu}
\label{sec:tg-reeb-def}

\begin{definicja}[Graf Reeba]
Niech $M$ będzie niepustą, zwartą, spójną rozmaitością gładką,
a $f:M\to\mathbb R$ funkcją gładką. Na $M$ definiujemy
$x\sim_f y$ wtedy i tylko wtedy, gdy $f(x)=f(y)$
oraz $x,y$ należą do tej samej spójnej składowej
zbioru poziomicowego $f^{-1}(f(x))$.
Przestrzeń ilorazową $R_f:=M/{\sim_f}$ nazywamy
\emph{przestrzenią Reeba} funkcji $f$. Odwzorowanie
ilorazowe oznaczamy $q_f:M\to R_f$.
Funkcja $f$ indukuje ciągłą funkcję
$\bar f:R_f\to\mathbb R$ spełniającą
$f=\bar f\circ q_f$.
Relacja jest równoważnością, ponieważ składowe spójne
każdej poziomicy tworzą jej rozbicie. Funkcja $\bar f$
jest jednoznacznie określona przez $\bar f([x])=f(x)$;
jest ciągła z definicji topologii ilorazowej.
\end{definicja}

Nie każda przestrzeń Reeba funkcji gładkiej jest
skończonym grafem. Najpierw przyjmiemy warunki,
które pozwalają udowodnić tę własność i opisać
stopnie jej wierzchołków.

\begin{definicja}[Prosta funkcja Morse'a na powierzchni]
Niech $M$ będzie zwartą gładką powierzchnią bez brzegu.
Gładką funkcję $f:M\to\mathbb R$ nazywamy \emph{funkcją Morse'a},
gdy jej wszystkie punkty krytyczne są niezdegenerowane.
Nazywamy ją \emph{prostą}, jeżeli dodatkowo różne punkty
krytyczne mają różne wartości. Indeks punktu krytycznego
jest liczbą ujemnych wartości własnych hesjanu w tym punkcie,
liczonych z krotnościami. Niezdegenerowanie oznacza
nieosobliwość hesjanu. W punkcie krytycznym zmiana mapy
przekształca macierz hesjanu przez kongruencję, więc
indeks jest niezależny od wyboru mapy na mocy prawa
bezwładności Sylvestera.
Na powierzchni wynosi on $0$, $1$ lub $2$.
\end{definicja}

\begin{stwierdzenie}[Regularne przedziały poziomów]
Niech $M$ będzie zwartą rozmaitością bez brzegu,
a $f:M\to\mathbb R$ gładką funkcją. Jeśli
$a<b$, przedział $[a,b]$ nie zawiera wartości krytycznych i
$f^{-1}([a,b])\ne\varnothing$, to istnieje
homeomorfizm, a nawet dyfeomorfizm
$f^{-1}([a,b])\cong f^{-1}(a)\times[a,b]$
zachowujący wartości funkcji $f$ jako drugą współrzędną.
\end{stwierdzenie}

\begin{proof}
Wybieramy metrykę riemannowską na $M$ i rozważamy
pole $X:=\nabla f/\|\nabla f\|^2$ na otoczeniu
$f^{-1}([a,b])$. Mianownik nigdzie tam nie znika.
Dla krzywej całkowej $\gamma$ pola $X$ mamy
$\frac{\mathrm d}{\mathrm dt}(f\circ\gamma)
=\mathrm df(X)=1$.
A zatem przepływ $\Phi_t$ przesuwa wartość funkcji
o dokładnie $t$. Aby uzasadnić istnienie tych odcinków przepływu,
wybieramy gładką funkcję odcięcia równą $1$ na
otoczeniu zwartej warstwy $f^{-1}([a,b])$, o nośniku
w obszarze, gdzie $\nabla f\ne0$. Mnożymy $X$ przez
tę funkcję i przedłużamy przez zero na $M$.
Otrzymane gładkie pole jest zupełne, ponieważ $M$
jest zwarta i bez brzegu. Dopóki wartość $f$ należy
do $[a,b]$, pochodna wzdłuż przepływu wynosi $1$.
Z jednoznaczności rozwiązań i ciągłości wynika,
że startując z poziomu $s\in[a,b]$, dla czasów
$t\in[a-s,b-s]$ pozostajemy w warstwie i mamy
$f(\Phi_t(x))=s+t$. W szczególności każdy punkt
niepustej warstwy można przesunąć na poziom $a$.
Odwzorowanie $(x,t)\mapsto\Phi_{t-a}(x)$,
$x\in f^{-1}(a)$, ma odwrotność
$y\mapsto(\Phi_{a-f(y)}(y),f(y))$.
Gładkość obu odwzorowań wynika z gładkiej zależności
przepływu od warunków początkowych.
\end{proof}

\begin{lemat}[Pas wokół pojedynczego punktu krytycznego]
\label{lem:tg-reeb-critical-band}
Niech $M$ będzie zwartą, orientowalną powierzchnią bez brzegu,
$f:M\to\mathbb R$ prostą funkcją Morse'a, a $p$ jej punktem
krytycznym o wartości $c$. Wybierzmy $\varepsilon>0$ tak, aby
$a=c-\varepsilon$ i $b=c+\varepsilon$ były wartościami regularnymi
i aby $[a,b]$ nie zawierał żadnej innej wartości krytycznej.
Składowa $W$ zbioru $f^{-1}([a,b])$ zawierająca $p$ jest:
\begin{enumerate}[label=(\roman*)]
\item dyskiem z jedną składową brzegu na poziomie $b$,
jeżeli $p$ jest minimum, albo z jedną składową na poziomie
$a$, jeżeli jest maksimum;
\item powierzchnią homeomorficzną z parą spodni, jeśli
$p$ jest siodłem; spośród trzech składowych jej brzegu
jedna leży na jednym z poziomów $a,b$, a dwie na drugim.
\end{enumerate}
Każda inna składowa tego pasa, jeżeli istnieje, jest
homeomorficzna z cylindrem $S^1\times[a,b]$.
\end{lemat}

\begin{proof}
Punkty brzegu pasa są regularne, zatem twierdzenie
o funkcji uwikłanej daje tam mapy półpłaszczyzny,
a wewnątrz pasa mamy zwykłe mapy powierzchni.
Pas jest zwartą powierzchnią z brzegiem, a jego
składowe są otwarte i domknięte w pasie, więc jest
ich skończenie wiele. Orientację dziedziczą z $M$.

Korzystamy z lematu Morse'a oraz klasycznego twierdzenia
o dołączaniu uchwytów w teorii Morse'a (np. Y. Matsumoto,
\emph{An Introduction to Morse Theory}, AMS, 2002).
W wersji dla zwartego pasa między regularnymi poziomami,
z jednym niezdegenerowanym punktem krytycznym indeksu
$\lambda$, twierdzenie opisuje pas jako dolny brzeg
razy przedział z dołączonym jednym uchwytem
$D^\lambda\times D^{2-\lambda}$ wzdłuż
$\partial D^\lambda\times D^{2-\lambda}$ i z
wygładzonymi narożnikami. Założenia są spełnione:
pas jest zwarty, $f$ nie ma punktów krytycznych na
brzegu i jedynym punktem krytycznym wewnątrz jest $p$.

Dla minimum uchwyt indeksu zero jest dyskiem
dołączanym wzdłuż zbioru pustego, więc stanowi osobną
składową z górnym brzegiem będącym okręgiem.
Dla maksimum stosujemy ten sam argument do $-f$.
Przy siodle uchwyt jest prostokątem przyczepionym
wzdłuż dwóch rozłącznych odcinków. Składowa $W$
składa się z jednego lub dwóch cylindrów dolnego
brzegu oraz tego uchwytu: pozostałe cylindry nie
spotykają uchwytu. Z addytywności charakterystyki
Eulera dla sklejenia wzdłuż podkompleksów (można
triangulować te powierzchnie zgodnie z odcinkami
sklejenia) otrzymujemy
\[
\chi(W)=0+\chi(D^1\times D^1)
-\chi(\partial D^1\times D^1)=0+1-2=-1.
\]

Ponieważ $p$ jest siodłem, w $W$ istnieją punkty
zarówno poniżej, jak i powyżej $c$. Z lokalnego modelu $-x^2+y^2$ lematu Morse'a
wybieramy w $W$ punkty o wartościach $s<c<t$.
Na regularnych przedziałach $[a,s]$ i $[t,b]$
poprzednie stwierdzenie daje odcinki przepływu
do poziomów $a$ i $b$; są zawarte w pasie, więc
pozostają w tej samej składowej $W$.
Niech $b_a,b_b\ge1$ oznaczają liczby okręgów
składających się na $W\cap f^{-1}(a)$ oraz
$W\cap f^{-1}(b)$. Wtedy brzeg $W$ ma
$r=b_a+b_b\ge2$ składowe.
Z orientowalności $W$ i klasycznego wzoru na
charakterystykę Eulera zwartej powierzchni
orientowalnej rodzaju $g$ z $r$ składowymi brzegu:
\[
-1=\chi(W)=2-2g-r,
\qquad 2g+r=3.
\]
Ponieważ $g\ge0$ jest całkowite, a $r\ge2$,
otrzymujemy $g=0$ i $r=3$.
Zatem $(b_a,b_b)=(1,2)$ albo $(2,1)$,
a klasyfikacja zwartych powierzchni z brzegiem
identyfikuje $W$ z parą spodni.
Wskazaliśmy przy tym dokładnie, gdzie
niezbędna jest orientowalność: dla powierzchni
nieorientowalnej dopuszczalne byłoby
$(b_a,b_b)=(1,1)$.

Każda składowa niezawierająca $p$ jest jednym
z cylindrów dolnego brzegu, do którego nie przyczepiono
uchwytu. W przypadku maksimum stosujemy opis
odwróconego pasa. Nie stosujemy twierdzenia o regularnych
przedziałach do całej warstwy zawierającej $p$.
\end{proof}

\begin{twierdzenie}[Skończoność i lokalna struktura grafu Reeba]
\label{thm:tg-reeb-finite}
Niech $M$ będzie zwartą, spójną, orientowalną
powierzchnią bez brzegu, a $f:M\to\mathbb R$
prostą funkcją Morse'a. Wtedy $R_f$ jest
homeomorficzna ze skończonym, spójnym grafem.
Przy kanonicznym wyborze wierzchołków nad punktami
krytycznymi wierzchołek odpowiadający minimum lub
maksimum ma stopień jeden, a odpowiadający punktowi
siodłowemu ma stopień trzy. Wnętrza krawędzi
odpowiadają rodzinom okręgów będących składowymi
poziomic niezawierającymi punktów krytycznych
(także na poziomach krytycznych innych składowych).
Dopuszczamy krawędzie równoległe.
\end{twierdzenie}

\begin{proof}
Punkty krytyczne funkcji Morse'a są izolowane,
a ich zbiór jest domknięty, gdyż jest zbiorem zer
różniczki $f$. Jest więc zwartą przestrzenią dyskretną,
a zatem skończoną.
Wybieramy regularne liczby
$a_0<a_1<\cdots<a_N$ rozdzielające kolejne wartości
krytyczne i rozbijamy powierzchnię na pasy zawierające
po jednej wartości krytycznej oraz pasy bez punktów
krytycznych. Na każdym zamkniętym przedziale regularnym
poprzednie stwierdzenie identyfikuje podzbiór powierzchni
z produktem poziomicy i przedziału. Regularna poziomica
na zwartej powierzchni bez brzegu jest skończoną sumą
okręgów. Po zidentyfikowaniu każdej składowej okręgu
z punktem produkt przechodzi w skończoną sumę odcinków.

Dla każdego przedziału $[a,b]$ obejmującego dokładnie jeden
punkt krytyczny $p$ o wartości $c$ stosujemy lemat
\ref{lem:tg-reeb-critical-band}. Sama topologia pasa
nie wystarcza jeszcze do wyznaczenia jego ilorazu
Reeba; sprawdzimy także włókno krytyczne.

Każda składowa poziomicy $f^{-1}(c)$ niezawierająca
$p$ jest zwartą regularną jednowymiarową rozmaitością
bez brzegu, a więc okręgiem. Istotnie, składowa jest
domknięta w zwartej poziomicy, a regularne mapy
poziomicy dają jej lokalne otwarte odcinki.
Taki okrąg $C$ można przesunąć przepływem na wszystkie
poziomy pasa. Przepływ istnieje jednolicie dla małych
czasów na otoczeniu zwartego $C$, a następnie
przedłużamy go w obie strony po regularnych
przedziałach aż do $a$ i $b$.
Otrzymany cylinder jest otwarty w pasie (lokalnie
jest produktem z przedziałem względnie otwartym
w $[a,b]$) i domknięty przez zwartość. Nie zawiera $p$,
więc nie może leżeć w składowej $W$ zawierającej $p$.
W konsekwencji $W\cap f^{-1}(c)$ jest spójne.

Po każdej stronie $c$ składowe poziomic w $W$ tworzą
rodziny okręgów, trywialne nad przedziałami $[a,c)$
i $(c,b]$: wybieramy poziom bazowy i używamy
przepływu na kolejnych zwartych podprzedziałach,
co daje zgodne identyfikacje. Jest ich odpowiednio
$b_a$ i $b_b$ (zero po jednej stronie ekstremum).
Tworzymy gwiazdę z $b_a+b_b$ ramion, po jednym
dla każdej rodziny, z parametrem wysokości $f$
na każdym ramieniu i wspólnym środkiem o wysokości $c$.
Odwzorowanie z $W$ do tej gwiazdy jest ciągłe poza
poziomem $c$ dzięki produktom, a na poziomie $c$
dzięki temu, że odległość od środka gwiazdy wynosi
$|f-c|$. Jego włóknami są dokładnie spójne składowe
poziomic w $W$. Indukuje zatem ciągłą bijekcję
z ilorazu Reeba pasa $W$ na gwiazdę, która jest
homeomorfizmem: dziedzina jest zwarta, a cel
Hausdorffa. Dla ekstremum otrzymujemy jedno ramię,
a dla siodła trzy. Pozostałe cylindry pasa dają
odcinki, ponieważ ich poziomice są spójne
i regularne; stosujemy przepływ na tych składowych.
To rozumowanie wykorzystuje orientowalność i rozdzielenie
wartości krytycznych: sam lokalny wzór hesjanu
bez kontroli sposobu sklejenia regularnych poziomic
nie wystarczyłby do wyznaczenia stopnia wierzchołka.
Ponieważ mamy skończenie wiele krytycznych poziomów,
po sklejeniu skończenie wielu odcinków i wierzchołków
uzyskujemy skończony graf. Dokładniej, końce
odcinków są utożsamiane wtedy i tylko wtedy, gdy
odpowiadają tej samej składowej wspólnej regularnej
poziomicy sąsiednich pasów. Otrzymujemy skończony
kompleks jednowymiarowy Hausdorffa. Lokalne mapy
do jego części są zgodne na wspólnych brzegach,
więc lemat o sklejaniu na skończonej rodzinie
domkniętych pasów daje ciągłą mapę z $M$.
Jej włókna są dokładnie klasami $\sim_f$, więc
indukuje homeomorfizm z $R_f$ na graf (ciągła
bijekcja ze zwartej przestrzeni do Hausdorffa).
Wierzchołki pomocnicze na regularnych poziomach
stopnia dwa można usunąć, łącząc sąsiednie odcinki.
Jego spójność wynika z
ciągłości i surjektywności odwzorowania ilorazowego $q_f$.
\end{proof}

\begin{samepage}
\begin{center}
\begin{minipage}{.96\linewidth}\centering
\begin{tikzpicture}[scale=1,>=Stealth,line cap=round]
 \begin{scope}[shift={(-2.6,0)}]
  \draw[tgEdge] (0,-1.25)--(0,1.25);
  \node[tgBlueVertex] at (0,-1.25) {};
  \node[tgRedVertex] at (0,1.25) {};
  \node[text=tgNavy,font=\small] at (0,-1.65) {sfera: odcinek};
 \end{scope}
 \begin{scope}[shift={(2.5,0)}]
  \node[tgBlueVertex] (b) at (0,-1.25) {};
  \node[tgTealVertex] (s1) at (0,-.35) {};
  \node[tgTealVertex] (s2) at (0,.65) {};
  \node[tgRedVertex] (t) at (0,1.4) {};
  \draw[tgEdge] (b)--(s1) (s2)--(t);
  \draw[tgBlue,line width=2pt] (s1) .. controls (-1.10,-.35) and (-1.10,.65) .. (s2);
  \draw[tgCoral,line width=2pt] (s1) .. controls (1.10,-.35) and (1.10,.65) .. (s2);
  \node[text=tgNavy,font=\small] at (0,-1.65) {torus: jeden cykl};
 \end{scope}
\end{tikzpicture}\par\smallskip

{\footnotesize Schemat grafów Reeba standardowych funkcji wysokości;
nie są to rysunki poziomic w samej powierzchni.}
\end{minipage}
\end{center}
\end{samepage}

\section{Rodzaj powierzchni zapisany w grafie Reeba}
\label{sec:tg-reeb-genus}

\begin{uwaga}[Znaczenie prostoty funkcji Morse'a]
W wynikach o grafach Reeba zakładamy, że każda wartość
krytyczna odpowiada dokładnie jednemu punktowi krytycznemu.
Bez tego założenia kilka lokalnych zmian zbioru poziomicowego
może zachodzić jednocześnie; graf Reeba nadal bywa skończony,
ale opis stopni jego wierzchołków wymaga wtedy dodatkowej
analizy. Wzoru na liczbę cykli poniżej nie należy przenosić
bez uzasadnienia na funkcje zdegenerowane ani na powierzchnie
z brzegiem.
\end{uwaga}

\begin{twierdzenie}[Pierwsza liczba Bettiego grafu Reeba]
\label{thm:tg-reeb-betti}
Niech $M=\Sigma_g$ będzie zwartą, spójną, orientowalną
powierzchnią bez brzegu rodzaju $g\ge0$, a
$f:M\to\mathbb R$ prostą funkcją Morse'a. Wtedy
\[
\dim H_1(R_f;\mathbb R)=g.
\]
W szczególności graf Reeba na sferze jest drzewem,
a na torusie ma dokładnie jeden niezależny cykl.
\end{twierdzenie}

\begin{proof}
Niech $c_i$ oznacza liczbę punktów krytycznych indeksu
$i\in\{0,1,2\}$. Z poprzedniego twierdzenia $R_f$ ma
$c_0+c_2$ wierzchołków stopnia jeden oraz $c_1$
wierzchołków stopnia trzy. Jeśli $v$ i $e$ oznaczają
liczbę jego wierzchołków i krawędzi, to
$v=c_0+c_1+c_2$ i
$2e=(c_0+c_2)+3c_1$.
Dla grafu spójnego otrzymujemy więc
\[
\dim H_1(R_f;\mathbb R)=e-v+1
=1+\frac{c_1-c_0-c_2}{2}.
\]
Z klasycznej teorii Morse'a dla powierzchni (albo
z rozkładu na uchwyty podpoziomic) wiemy, że
$\chi(\Sigma_g)=c_0-c_1+c_2=2-2g$:
minimum dodaje dysk, siodło uchwyt jednowymiarowy,
a maksimum dysk zamykający składową brzegu, co zmienia
charakterystykę Eulera odpowiednio o $+1,-1,+1$.
Zatem $c_1-c_0-c_2=2g-2$. Po podstawieniu uzyskujemy
$\dim H_1(R_f;\mathbb R)=g$, jak należało dowieść.
\end{proof}

\begin{przyklad}[Sfera i torus]
Standardowa funkcja wysokości na sferze ma jedno minimum
i jedno maksimum, a więc $(c_0,c_1,c_2)=(1,0,1)$.
Graf Reeba ma dwa wierzchołki i jedną krawędź,
zatem $\beta_1=0$. Dla torusa ustawionego pionowo
standardowa funkcja wysokości ma jedno minimum,
dwa siodła i jedno maksimum, czyli $(1,2,1)$.
Wtedy $v=4$, $e=4$ i $\beta_1=4-4+1=1$.
Dodawanie par krytycznych punktów podczas deformacji
funkcji może zwiększyć liczby $c_i$, ale równość
$c_0-c_1+c_2=2-2g$ i wartość $\beta_1=g$ pozostają.
\end{przyklad}

\begin{uwaga}[Zakres twierdzenia]
Ostatnia równość jest szczególną własnością prostych
funkcji Morse'a na \emph{orientowalnych powierzchniach
zamkniętych}. Dla rozmaitości wyższego wymiaru
lub funkcji o zdegenerowanych krytycznych poziomach
graf Reeba może mieć inną kombinatorykę, a powyższe
liczenie stopni wierzchołków nie ma zastosowania.
\end{uwaga}

\begin{przyklad}[Dlaczego potrzebna jest orientowalność]
Na płaszczyźnie rzutowej $\mathbb{RP}^2$ istnieje funkcja
Morse'a z dokładnie jednym minimum, jednym siodłem
i jednym maksimum; odpowiada to standardowemu rozkładowi
komórkowemu z jedną komórką w każdym wymiarze
$0,1,2$. Przekroczenie siodła dołącza w tym przypadku
uchwyt skręcony: jedna składowa regularnej poziomicy
pozostaje jedną składową. Graf Reeba jest odcinkiem,
a obraz siodła ma stopień dwa, nie trzy.
Rzeczywiście $c_0-c_1+c_2=1=\chi(\mathbb{RP}^2)$;
wzór z Twierdzenia~\ref{thm:tg-reeb-betti} nie
może zostać zastosowany bez orientowalności.
\end{przyklad}

\section{Graf Reeba a rozkład na uchwyty}
\label{sec:tg-reeb-handles}

\begin{przyklad}[Siodło łączy lub rozdziela]
Niech $a<b$ będą regularnymi wartościami funkcji
Morse'a $f:\Sigma_g\to\mathbb R$ i niech przedział
$(a,b)$ zawiera dokładnie jeden punkt krytyczny indeksu
jeden. Zbiór podpoziomowy $f^{-1}((-\infty,a])$ po przekroczeniu
siodła otrzymuje jeden uchwyt jednowymiarowy
$D^1\times D^1$, przyczepiony wzdłuż
$\partial D^1\times D^1$. Jeżeli jego końce są przyczepione
do dwóch różnych składowych brzegowych podpoziomicy,
dwie składowe regularnego poziomu łączą się w jedną.
Jeżeli końce leżą na tej samej składowej brzegu,
powstają dwie składowe poziomu. Obie operacje mają
w grafie Reeba lokalny obraz wierzchołka stopnia trzy.
\end{przyklad}

\begin{samepage}
\begin{center}
\begin{minipage}{.96\linewidth}\centering
\begin{tikzpicture}[scale=.98,>=Stealth,line cap=round,line join=round]
 \begin{scope}[shift={(-2.6,0)}]
  \draw[tgBlue,line width=2pt] (-1,-.8) .. controls (-.8,-.1) and (-.6,.25) .. (0,1);
  \draw[tgTeal,line width=2pt] (1,-.8) .. controls (.8,-.1) and (.6,.25) .. (0,1);
  \draw[tgEdge] (0,1)--(0,1.65);
  \node[tgRedVertex] at (0,1) {};
  \node[text=tgNavy,font=\small] at (0,-1.85) {połączenie: $2\to1$};
 \end{scope}
 \begin{scope}[shift={(2.6,0)}]
  \draw[tgBlue,line width=2pt] (0,-.8) .. controls (-.6,-.05) and (-.8,.45) .. (-1,1);
  \draw[tgTeal,line width=2pt] (0,-.8) .. controls (.6,-.05) and (.8,.45) .. (1,1);
  \draw[tgEdge] (0,-1.45)--(0,-.8);
  \node[tgRedVertex] at (0,-.8) {};
  \node[text=tgNavy,font=\small] at (0,-1.85) {rozdzielenie: $1\to2$};
 \end{scope}
\end{tikzpicture}\par\smallskip

{\footnotesize Lokalny obraz wierzchołka siodłowego
w grafie Reeba; wartości funkcji rosną od dołu ku górze.}
\end{minipage}
\end{center}
\end{samepage}

\begin{uwaga}[Czego graf Reeba nie zachowuje]
Iloraz $R_f$ zapamiętuje liczbę i sposób łączenia
spójnych składowych poziomic, ale nie pełną geometrię
metryczną powierzchni. Dwie różne funkcje Morse'a
mogą mieć izomorficzne grafy Reeba; dodatkowe dane,
na przykład wartości krytyczne i rodzaje składowych
podpoziomic, pozwalają na dokładniejszą klasyfikację.
Na powierzchni orientowalnej sam pierwszy wymiar homologii
grafu wyznacza rodzaj powierzchni na mocy
Twierdzenia~\ref{thm:tg-reeb-betti}.
\end{uwaga}

